\mychapter{Wprowadzenie}
\label{chap:intro}

\section{Miasto $t$ minut: kontekst i motywacja}
\label{sec:intro:kontekst}

Koncepcja \emph{miasta 15 minut} (częściej uogólniana jako \emph{miasto $t$ minut})
opisuje wizję organizacji przestrzeni miejskiej, w której podstawowe potrzeby życia codziennego
--- takie jak zakupy, edukacja, opieka zdrowotna, usługi publiczne, rekreacja czy gastronomia ---
mogą być zaspokojone w krótkim czasie podróży, zazwyczaj pieszo lub rowerem.
Idea ta została spopularyzowana m.in.\ przez Moreno w 2016 roku \cite{Moreno2016QuartHeure}
i od tego czasu stała się jednym z częściej dyskutowanych paradygmatów w urbanistyce,
politykach miejskich oraz analizach dostępności.

Z punktu widzenia praktyki planistycznej koncepcja miasta $t$ minut bywa sprowadzana do
pytania, które ma bezpośrednią interpretację obliczeniową:
\begin{quote}
\emph{Z jakich miejsc w mieście da się dotrzeć do zadanych kategorii usług w czasie
nieprzekraczającym $t$ minut?}
\end{quote}
Odpowiedź na to pytanie jest istotna zarówno diagnostycznie (pomiar dostępności usług i jej nierówności),
jak i operacyjnie (wskazanie miejsc, w których potrzebne są interwencje: dogęszczenie usług,
zmiana organizacji ruchu, poprawa spójności sieci pieszej, likwidacja barier itp.).

W literaturze spotyka się różne warianty problemu miasta $t$ minut.
Różnią się one m.in.\ zestawem kategorii usług, sposobem modelowania czasu przemieszczania (pieszo/rowerem)
oraz sposobem agregacji wyniku do postaci wskaźników porównawczych lub map tematycznych.
Niezależnie jednak od wariantu, rdzeniem problemu pozostaje wyznaczanie zasięgów czasowych
(\emph{izochron}) w grafie reprezentującym miasto.

\clearpage

Z perspektywy informatycznej miasto $t$ minut jest interesujące z co najmniej dwóch powodów.
Po pierwsze, stanowi przykład zadania, w którym klasyczne problemy grafowe (głównie problem najkrótszej ścieżki) przekładają się bezpośrednio na interpretację społeczną i przestrzenną.
Po drugie, praktyczne zastosowania wymagają obliczeń na dużych grafach miejskich oraz wykonywania
wielu podobnych zapytań (dla wielu kategorii usług i wielu progów czasu),
co rodzi pytania o efektywność algorytmów i implementacji.

\section{Model obliczeniowy i formalizacja problemu}
\label{sec:intro:formalizacja}

W formalizacji algorytmicznej kluczowe jest rozróżnienie pomiędzy odległością euklidesową
(„w linii prostej”) a odległością \emph{sieciową}, wynikającą z realnego przebiegu ulic, przejść
i skrzyżowań. Naturalnym językiem opisu takiej sieci jest teoria grafów \cite{barbieri2023graph}.
W niniejszej pracy miasto modelowane jest jako graf ważony
$G=(V,E)$, gdzie wierzchołki $V$ odpowiadają punktom sieci (np.\ skrzyżowaniom, wejściom do przejść,
węzłom ciągów pieszych), a krawędzie $E$ odpowiadają odcinkom możliwym do pokonania.
Waga $w(e)$ krawędzi $e\in E$ jest interpretowana jako koszt przejścia, najczęściej czas w minutach,
wyznaczany na podstawie długości odcinka oraz przyjętej prędkości przemieszczania się.

Niech $S_k\subseteq V$ oznacza zbiór lokalizacji usług danej kategorii $k$
(np.\ aptek, sklepów spożywczych, szkół). Dla wierzchołka $v\in V$ definiujemy czas dostępu do kategorii $k$ jako
\[
\tau_k(v)=\min_{s\in S_k} d_G(v,s),
\]
gdzie $d_G(v,s)$ jest najkrótszą odległością w grafie $G$ według wag czasowych.
Wtedy zbiór wierzchołków spełniających warunek dostępności kategorii $k$ w czasie $t$ ma postać
\[
C_k(t)=\{\,v\in V:\ \tau_k(v)\le t\,\}.
\]
Jeżeli interesuje nas jednoczesna dostępność wielu kategorii usług $k=1,\dots,K$,
to definiujemy
\[
C(t)=\bigcap_{k=1}^{K} C_k(t),
\]
czyli zbiór wierzchołków, z których da się dotrzeć do \emph{każdej} wymaganej kategorii usług w czasie $\le t$.

\clearpage

W literaturze spotyka się również interpretację, w której zamiast rozpatrywać tylko najbliższą usługę,
analizuje się zasięgi czasowe wokół \emph{konkretnych lokalizacji} usług.
To podejście jest szczególnie wygodne obliczeniowo w sytuacji, gdy liczba usług jest mniejsza
niż liczba potencjalnych punktów startowych (np.\ adresów lub węzłów sieci) \cite{barbieri2023graph}.
Niech $C^k_s(t)$ oznacza zbiór wierzchołków, z których można dotrzeć do konkretnej usługi $s\in S_k$ w czasie $\le t$:
\[
C^k_s(t)=\{\,v\in V:\ d_G(v,s)\le t\,\}.
\]
Wówczas mamy równoważność:
\[
C_k(t)=\bigcup_{s\in S_k} C^k_s(t),
\]
a więc wynik dla kategorii jest sumą zasięgów wokół jej lokalizacji.
Takie odwrócenie perspektywy (start od usług zamiast od badanych wierzchołków) jest ważne praktycznie,
ponieważ umożliwia budowanie zbioru $C(t)$ przez serię obliczeń izochron.

Wyniki analiz miasta $t$ minut bywają następnie agregowane do postaci wskaźników porównawczych
lub map tematycznych. Przykładem są podejścia oparte na grafie indukowanym przez wierzchołki $C(t)$
i analizie jego własności (np.\ rozłącznych składowych jako „luk” w dostępności) \cite{barbieri2023graph},
a także metody siatkowe (np.\ heksagonalna tessalacja) i indeksy typu NEXI,
w których czasy dostępu są uśredniane w komórkach siatki \cite{olivari2023presenting}.

\section{Wyzwania algorytmiczne w analizach miasta $t$ minut}
\label{sec:intro:wyzwania}

Choć definicje zbiorów $C_k(t)$ i $C(t)$ są proste, ich wyznaczanie na realnych danych miejskich
wiąże się z istotnymi wyzwaniami obliczeniowymi.

\subsection{Lokalność izochron i poszukiwanie algorytmów „dopasowanych”}

W analizach miasta $t$ minut interesuje nas przede wszystkim informacja o osiągalności
\emph{w zadanym progu czasu} $t$, a nie pełna tablica odległości do wszystkich wierzchołków.
Dla typowych wartości $t$ (np.\ 10--20 minut marszu) izochrona obejmuje relatywnie niewielki fragment miasta.
Ta własność \emph{lokalności} sugeruje, że algorytmy i implementacje powinny wykorzystywać fakt,
iż większość grafu nie jest odwiedzana w trakcie obliczeń dla danego źródła.

Klasycznym narzędziem do wyznaczania najkrótszych odległości w grafach z nieujemnymi wagami
jest algorytm Dijkstry \cite{dijkstra}. W kontekście miasta $t$ minut często stosuje się jego warianty,
w których przeszukiwanie jest ograniczane progiem czasu (tzw.\ \emph{$t$-bounded Dijkstra})
oraz warianty, które minimalizują narzut pamięciowy poprzez wstawianie do kolejki priorytetowej
tylko wierzchołków faktycznie odkrytych (podejście znane jako \emph{Uniform Cost Search})
\cite{Lam2024}. W literaturze pojawiają się również triki wieloźródłowe,
w których wiele lokalizacji usług w obrębie jednej kategorii jest obsługiwanych przez jedno uruchomienie
algorytmu (np.\ przez dodanie superźródła połączonego krawędziami o koszcie $0$)
\cite{Lam2024}.

Jednocześnie, postęp w projektowaniu algorytmów najkrótszych ścieżek oraz metod przeszukiwania grafów
powoduje, że pojawiają się nowe propozycje, które mogą być lepiej dopasowane do specyfiki problemu miasta $t$ minut
niż klasyczne podejście oparte wyłącznie na Dijkstrze. W szczególności interesujące są algorytmy,
które w pewnych klasach grafów lub przy pewnym typie zapytań oferują lepszą skalowalność
czasową lub pamięciową.
W niniejszej pracy rozpatrywany jest nowy algorytm zaproponowany w pracy \cite{arxiv2504_17033},
który zostanie zestawiony z podejściem opartym na Dijkstrze w zadaniu wyznaczania obszarów miasta $t$ minut.

\subsection{Co oznacza „lepszy algorytm” w praktyce miasta $t$ minut?}

W kontekście analizy miasta $t$ minut samo porównanie asymptotycznej złożoności algorytmów bywa niewystarczające.
O ostatecznej przydatności metody decydują m.in.:
\begin{itemize}
  \item \textbf{czas całkowity} wyznaczenia $C(t)$ dla pełnego zestawu kategorii usług,
  \item \textbf{zużycie pamięci} (szczególnie istotne przy dużych grafach),
  \item \textbf{skalowalność} względem liczby usług, kategorii i rozmiaru grafu,
  \item \textbf{zgodność wyniku} (czy algorytmy wyznaczają identyczne zbiory $C(t)$ przy tych samych założeniach).
\end{itemize}
Z tych powodów praca koncentruje się na porównaniu algorytmów w warunkach możliwie zbliżonych do praktycznego użycia
w analizach dostępności, a nie jedynie na analizie teoretycznej.

\subsection{Wpływ implementacji: Python a Rust}
\label{sec:intro:py-vs-rust}

Wybór optymalnego podejścia technicznego do problemu najktótszej ścieżki nie sprowadza się wyłącznie do złożoności obliczeniowej,
ale w praktyce jest silnie zależny od doboru języka programowania, struktur danych oraz wykonawczego środowiska
uruchomieniowego. Z tego względu w niniejszej pracy rozważane będzie nie tylko porównanie algorytmów na poziomie idei,
lecz także porównanie ich implementacji w dwóch odmiennych ekosystemach: Python oraz Rust.

W pracy Lama \cite{Lam2024} autor porównał implementacje tego samego podejścia w Pythonie i w Ruście,
uzyskując w swoich eksperymentach kilkukrotne przyspieszenie wersji w Rust (około trzykrotne względem wersji w Pythonie).
Wynik ten sugeruje, że w zadaniach obliczeniowo intensywnych, takich jak wielokrotne wyznaczanie najkrótszych ścieżek na dużych grafach miejskich,
wybór środowiska uruchomieniowego może istotnie wpływać na czas działania.
W niniejszej pracy aspekt ten zostanie przebadany również w kontekście alternatywnego algorytmu najkrótszej ścieżki:
celem jest ocena, czy obserwowane w literaturze przewagi implementacji niskopoziomowych (np.\ w Rust) utrzymują się
także dla innego podejścia algorytmicznego.

W ramach tej pracy planowane jest zbadanie, w jakim stopniu wybór języka (Python vs Rust) wpływa na:
(i) czas wyznaczania zbiorów $C_k(t)$,
(ii) zużycie pamięci

\subsection{Równoległość obliczeń w zadaniu miasta $t$ minut}
\label{sec:intro:parallel}

Kolejnym istotnym wymiarem praktycznej efektywności jest możliwość wykorzystania współczesnych procesorów wielordzeniowych.
W problemie miasta $t$ minut naturalną osią równoległości jest fakt, że obliczenia dla różnych typów usług
(sklepy, szkoły, opieka zdrowotna itd.) są wykonywane niezależnie, a dopiero na końcu łączone
w wynik zbioru $C(t)$.

W literaturze spotyka się wprost sugestię równoleglenia tego etapu: dla różnych typów usług można uruchamiać przeszukiwania
w grafie w oddzielnych wątkach, a następnie scalić uzyskane etykiety osiągalności. W pracy Lama wskazano, że bieżąca
implementacja jest jednordzeniowa, ale algorytm można zrównoleglić tak, aby różne typy usług były przeszukiwane przez
różne wątki; po zakończeniu każdy wątek zwraca tablicę etykiet $v.r$, a wątek główny łączy wyniki w rezultat końcowy
\cite{Lam2024}.

W niniejszej pracy planowane jest przeprowadzenie eksperymentów obejmujących:
\begin{itemize}
  \item implementację wielowątkową (z podziałem obliczeń po typach usług),
  \item pomiar przyspieszenia w funkcji liczby wątków,
  \item porównanie efektów równoległości w Python i w Rust.
\end{itemize}
Takie ujęcie pozwoli ocenić, czy z punktu widzenia analiz miejskich bardziej opłaca się inwestować w nowe algorytmy,
czy w optymalizację implementacji i równoległość obliczeń.

\section{Cel, zakres i pytania badawcze}
\label{sec:intro:cel}

\subsection{Cel główny}

Celem niniejszej pracy jest analiza i porównanie podejścia opartego na algorytmie Dijkstry
z nowym algorytmem zaproponowanym w \cite{arxiv2504_17033} w kontekście
wyznaczania obszarów miasta $t$ minut.

Porównanie jest prowadzone z perspektywy zadania obliczeniowego opisanego w
rozdz.\ \ref{sec:intro:formalizacja}, tj.\ wyznaczania zbiorów $C_k(t)$
dla zadanych kategorii usług i progu czasu $t$.

\subsection{Zakres pracy}

Zakres pracy obejmuje:
\begin{itemize}
  \item implementację podejścia bazowego opartego na Dijkstrze,
  \item implementację algorytmu z \cite{arxiv2504_17033}
        w wersji umożliwiającej rozwiązanie zadania miasta $t$ minut,
  \item zaprojektowanie i przeprowadzenie eksperymentów porównawczych,
  \item analizę wyników (czas, pamięć, skalowalność) i sformułowanie wniosków.
\end{itemize}
Poza zakresem pracy pozostają zagadnienia stricte urbanistyczne (np.\ dobór polityk miejskich),
a także ocena jakości danych źródłowych w sensie geograficznym; dane są traktowane jako wejście do problemu
algorytmicznego.

\subsection{Pytania badawcze}

Praca jest ukierunkowana na następujące pytania badawcze:
\begin{description}
  \item[P1.] Czy algorytm z \cite{arxiv2504_17033} pozwala obniżyć koszt obliczeń (czas i pamięć)
            w wyznaczaniu zbiorów $C_k(t)$ w problemie miasta $t$ minut
            w porównaniu do podejścia bazowego opartego na algorytmie Dijkstry,
            przy zachowaniu zgodności wyniku?

  \item[P2.] Jakie przyspieszenie dają równoległe obliczenia po typach usług (wielowątkowość)
            w wyznaczaniu $C_k(t)$ oraz jak skaluje się ono wraz z liczbą wątków
            w porównaniu do wersji jednowątkowej?

  \item[P3.] W jakim stopniu wybór języka implementacji (Python vs Rust) wpływa na czas działania,
            zużycie pamięci oraz skalowalność rozważanych rozwiązań w dużych grafach miejskich,
            i czy obserwowane różnice są porównywalne dla podejścia opartego na Dijkstrze
            oraz dla algorytmu z \cite{arxiv2504_17033}?
\end{description}

\clearpage

\section{Metodyka badań i założenia eksperymentalne}
\label{sec:intro:metodyka}

W celu udzielenia odpowiedzi na pytania badawcze przyjęto metodykę opartą na powtarzalnych eksperymentach
na grafach miejskich oraz zdefiniowanych zestawach usług. Eksperymenty są projektowane tak, aby porównywać algorytmy
w warunkach możliwie zbliżonych do praktycznych analiz miasta $t$ minut, przy zachowaniu porównywalności danych wejściowych
oraz parametrów modelu.

\subsection{Dane wejściowe i konstrukcja grafu}
\label{sec:intro:metodyka:dane}

Eksperymenty opierają się na grafach reprezentujących \emph{sieć pieszą} miasta, zbudowanych na podstawie danych
OpenStreetMap (OSM). Dane OSM są powszechnie wykorzystywane w analizach dostępności, ponieważ zawierają zarówno geometrię
sieci transportowej, jak i informacje o punktach zainteresowania (PoI), które mogą pełnić rolę usług w problemie miasta
$t$ minut. W niniejszej pracy ruch modelowany jest wyłącznie jako pieszy.

Na podstawie danych OSM konstruowany jest graf ważony $G=(V,E)$, w którym wierzchołki odpowiadają punktom sieci,
a krawędzie odcinkom możliwym do pokonania pieszo.
Wagi krawędzi są interpretowane jako czasy przejścia, wyznaczane na podstawie długości odcinka oraz przyjętej prędkości marszu.
Lokalizacje usług są wyznaczane na podstawie danych OSM i mapowane na elementy sieci w sposób zapewniający spójność
między porównywanymi algorytmami oraz implementacjami. Proces budowy grafu jest inspirowany podejściem zastosowanym przez
Lama \cite{Lam2024} i obejmuje kroki porządkujące topologię sieci oraz przygotowanie zbiorów węzłów odpowiadających usługom,
tak aby stanowiły one źródła/cele w obliczeniach izochron.

Szczegółowy opis procedury pozyskania danych, filtracji sieci do ruchu pieszego, konstrukcji wag czasowych oraz mapowania usług
do grafu zostanie przedstawiony w~rozdziale~\ref{chap:dane_i_model}.

\clearpage

\subsection{Zakres przestrzenny eksperymentów}
\label{sec:intro:metodyka:obszary}

Eksperymenty zostaną przeprowadzone dla trzech obszarów o zróżnicowanej skali i charakterystyce,
co pozwoli ocenić zachowanie badanych algorytmów zarówno w małych, lokalnych analizach,
jak i w warunkach dużego obciążenia obliczeniowego typowego dla metropolii.

\begin{itemize}
  \item \textbf{Gdańsk Południe} --- część Gdańska położona na południowym zachodzie miasta,
        obejmująca od 2018~r.\ cztery jednostki pomocnicze:
        Chełm, Jasień, Orunię Górną--Gdańsk Południe oraz Ujeścisko-Łostowice \cite{wiki:gdansk_poludnie}.
        Stanowi przykład relatywnie małego obszaru i~służy jako scenariusz „lokalny”,
        istotny z~punktu widzenia analiz planistycznych na poziomie dzielnicy.

  \item \textbf{Gdańsk (całe miasto)} --- scenariusz średniej skali, pozwalający badać wpływ wzrostu rozmiaru grafu
        oraz liczby usług na czas i zużycie pamięci, przy zachowaniu spójnego kontekstu urbanistycznego
        i porównywalnych założeń modelu.

  \item \textbf{Londyn} --- przykład bardzo dużego miasta, służący do oceny skalowalności podejść
        w warunkach grafu o znacznie większym rozmiarze oraz większej liczbie lokalizacji usług.
        Ten przypadek pełni rolę testu „stresowego” dla algorytmów oraz implementacji (w tym wariantów wielowątkowych).
\end{itemize}

Dokładny sposób wyznaczenia granic obszarów oraz pozyskania odpowiadających im danych z~OpenStreetMap
zostanie opisany w~rozdziale~\ref{chap:dane_i_model}.

\subsection{Projekt eksperymentów}
\label{sec:intro:metodyka:scenariusze}

W niniejszej pracy celem eksperymentów jest przede wszystkim porównanie \emph{własności obliczeniowych} rozważanych podejść,
a nie analiza wrażliwości wyników na parametry urbanistyczne lub modelowe.
Dlatego przyjęto stały próg czasu $t=15$ minut, ustaloną prędkość marszu $v=5\,\mathrm{km/h}$ (zgodnie z \cite{olivari2023presenting})
oraz jeden zdefiniowany zbiór kategorii usług.
Takie ograniczenie pozwala skupić się na porównaniu algorytmów i implementacji w warunkach odpowiadających typowemu użyciu
pojęcia „miasta 15 minut”.

Eksperymenty są projektowane jako zestaw porównań kontrolowanych, w których zmieniane są wyłącznie czynniki istotne
dla celów pracy:
\begin{itemize}
  \item \textbf{wariant algorytmiczny} --- podejście bazowe oparte na algorytmie Dijkstry oraz algorytm z \cite{arxiv2504_17033},
  \item \textbf{wariant implementacyjny} --- implementacja w Python oraz implementacja w Rust,
  \item \textbf{wariant wykonawczy} --- wykonanie sekwencyjne (jednowątkowe) oraz wykonanie równoległe (wielowątkowe),
        zgodnie z podziałem obliczeń po typach usług.
\end{itemize}
Skalowalność rozwiązań jest oceniana poprzez porównania na obszarach o różnej wielkości i złożoności sieci
(opisanych w podrozdziale \ref{sec:intro:metodyka:obszary}), co pozwala analizować wpływ rozmiaru grafu i liczby usług
bez zmiany założeń modelowych.

\subsection{Miary porównawcze}
\label{sec:intro:metodyka:miary}

W pracy raportowane są miary bezpośrednio związane z efektywnością obliczeniową:
\begin{itemize}
  \item \textbf{czas wykonania} wyznaczenia zbiorów osiągalności (zarówno w ujęciu całkowitym, jak i w ujęciu per kategoria),
  \item \textbf{zużycie pamięci} w trakcie obliczeń (np.\ wartość maksymalna oraz wartości reprezentatywne dla przebiegu),
  \item \textbf{skalowalność} czasu i pamięci w funkcji rozmiaru grafu oraz liczby usług (na przykładzie trzech obszarów testowych),
  \item \textbf{przyspieszenie równoległe} i efektywność w funkcji liczby wątków w wariancie wielowątkowym,
  \item \textbf{zgodność wyniku} pomiędzy podejściami (weryfikacja, że wyznaczane zbiory osiągalności są identyczne
        przy tych samych danych wejściowych i parametrach modelu).
\end{itemize}
Miary przestrzenne (np.\ udział węzłów spełniających warunek, miary pokrycia lub agregacja do siatki) mogą być wykorzystywane
pomocniczo do prezentacji i interpretacji rezultatów, jednak nie stanowią głównych kryteriów porównania w tej pracy.

\section{Struktura pracy}
\label{sec:intro:struktura}

Praca składa się z siedmiu rozdziałów.

W rozdziale~1 wprowadzono problem miasta 15 minut jako zadanie obliczeniowe, przedstawiono jego formalizację w modelu grafowym
oraz omówiono kluczowe wyzwania algorytmiczne. Zdefiniowano również cel pracy i pytania badawcze, a także przyjęte założenia
eksperymentalne.

Rozdział~2 opisuje dane i model obliczeniowy wykorzystywany w eksperymentach: sposób pozyskania i przygotowania sieci pieszej
oraz punktów usług, przyjęty model wag czasowych oraz elementy zapewniające
porównywalność eksperymentów pomiędzy różnymi algorytmami i implementacjami.

W rozdziale~3 przedstawiono podejście bazowe oparte na algorytmie Dijkstry w kontekście wyznaczania miasta 15 minut.
Omówiono warianty istotne praktycznie dla obliczania izochron oraz rozważono wersję sekwencyjną
i wariant równoległy oparty na podziale obliczeń po typach usług.

Rozdział~4 poświęcono algorytmowi zaproponowanemu w \cite{arxiv2504_17033}. Przedstawiono jego ideę, elementy kluczowe dla
wydajności oraz sposób adaptacji do zadania obliczania zbiorów osiągalności w problemie miasta 15 minut, a także zestawiono go
koncepcyjnie z podejściem bazowym.

W rozdziale~5 opisano szczegóły implementacji i środowisko eksperymentalne: architekturę rozwiązania, implementacje w Pythonie i w Ruście,
warianty jednowątkowe i wielowątkowe, procedury walidacji poprawności oraz metody pomiaru czasu i zużycia pamięci.

Rozdział~6 zawiera wyniki eksperymentów dla trzech obszarów o rosnącej skali (Gdańsk Południe, Gdańsk, Londyn).
Zaprezentowano porównanie algorytmów (Dijkstra vs \cite{arxiv2504_17033}), porównanie implementacji (Python vs Rust),
oraz wpływ równoległości obliczeń (jednowątkowość vs wielowątkowość), wraz z dyskusją uzyskanych rezultatów.

Pracę zamyka rozdział~7, w którym podsumowano wyniki, udzielono odpowiedzi na pytania badawcze, wskazano ograniczenia oraz
zaproponowano możliwe kierunki dalszych badań.

